\chapter{Introduction}
\label{cap:cap1}

% Introducerea va avea aproximativ 2 (două) pagini care vor conține motivația alegerii temei, relevanța și contextul temei alese, obiectivele generale ale lucrării, metodologia și instrumentele utilizate și o scurtă descriere a structurii lucrării (titlul capitolelor, o scurtă descriere și legătura dintre acestea).

% În acest capitol nu se introduc figuri, table sau listing-uri de cod. Pot fi în schimb referite!

% %\sout{text de șters}

% \section{Recomandări generale pentru redactare}

% \begin{description}
%     \item[Font:]  Times New Roman, mărimea 12
%     \item[Spațiere:] 1 între rânduri
%     \item[Margini:] 2 cm (sus, jos, stânga, dreapta)
%     \item[Număr de pagini:] Maxim 30 pagini fără rezumat, bibliografie și anexe
%     \item[Stil bibliografic:] IEEE
%     \item[Format:] Word sau LaTeX (în funcție de preferințe)
% \end{description}

% \section{Structura lucrării scrise}

% \begin{itemize}
%     \item Rezumatul lucrării (max. 1 pg.)
%     \item 1. Introducere (aprox. 2 pg.)
%     \item 2. Fundamentarea teoretică și soluții similare (aprox. 5 pg.)
%     \item 3. Soluția propusă (aprox 14 pg.)
%     \item 4. Testarea soluției și rezultate experimentale (aprox. 7 pg.)
%     \item 5. Concluzii (max. 2 pg.)
%     \item Bibliografie
%     \item Anexe
% \end{itemize}

% \section{Recomandări conținut introducere}

% \begin{itemize}
%     \item Contextul și importanța temei alese.
%     \item Obiectivele generale și specifice ale lucrării.
%     \item Metodologia de cercetare utilizată.
%     \item Structura generală a lucrării.
% \end{itemize}

\section{Context and Motivation}

Biomass estimation represents an important component of precision agriculture, being used for monitoring crop and pasture status, assessing productivity, and optimizing agricultural resource management processes \cite{castro2020deep, karila2022grass}. Traditional methods of biomass measurement are mainly based on harvesting and weighing plant material, procedures that are costly, time-consuming, and difficult to apply at large scales.

Recent advances in the fields of artificial intelligence, computer vision, and machine learning have opened up the possibility of estimating biomass directly from images. The development of models based on deep neural networks and the use of pre-trained visual representations have enabled promising results in numerous agricultural applications. Among the various available data sources, RGB images are particularly attractive due to their low cost, accessibility, and ease of acquisition under real-world operating conditions.

However, the high performance reported in the literature is, in most cases, evaluated using standard validation procedures, such as random cross-validation. In the presence of the temporal and spatial dependencies specific to agricultural data, these procedures can lead to overly optimistic estimates of actual performance \cite{roberts2017cross, ploton2020spatial}. Thus, a model that appears performant during validation may register significantly poorer results when applied to data collected in other periods or other locations.

From this perspective, the problem is no longer just building a model capable of estimating biomass, but also correctly assessing its generalization ability. In the absence of adequate validation protocols, the reported performance may provide a misleading picture of the model's behavior under real-world usage conditions.

\section{Relevance of the Topic}

The growing interest in using artificial intelligence in agriculture has led to the emergence of an increasing number of models dedicated to biomass estimation. However, most research focuses on improving predictive performance, paying little attention to how this performance is evaluated.

In practice, agricultural datasets are characterized by significant temporal and spatial variations. Weather conditions, soil characteristics, vegetation types, and the timing of image collection can introduce distribution shifts between the data used for training and those encountered later in deployment. In these situations, conventional validation methods may fail to provide a realistic estimate of future performance.

Therefore, the study of the reliability of validation protocols represents a relevant issue from both an academic and a practical perspective. The results of such an analysis can contribute to choosing more robust evaluation procedures and to a more accurate interpretation of the performance reported by machine learning models used in the agricultural domain.

\section{Research Objectives}

The main objective of this work is to analyze the impact of different validation protocols on the performance estimation of biomass prediction models under conditions of temporal and spatial data distribution shift.

In this context, temporal distribution shift refers to changes in the relationship between image appearance and biomass that occur over time. For example, pasture images collected in different seasons or years may differ in illumination, vegetation color, senescence stage, and species composition, even if they belong to similar plots. Spatial distribution shift refers to differences across geographic locations, such as soil type, climate, and pasture composition. In this thesis, spatial units are defined in terms of Australian states, since the dataset provides state-level metadata and the Date-State validation protocol uses state as a geographic grouping criterion. When a model is trained in one region or period and evaluated in another, its performance can degrade because the training data do not fully represent the conditions encountered in deployment.

The experimental context of this work is the CSIRO Image2Biomass Kaggle competition \cite{csiro-biomass}, organized by the Commonwealth Scientific and Industrial Research Organisation (CSIRO) of Australia. The objective of the competition is to estimate multiple pasture biomass components from top-view images of field quadrats. Participants are provided with a public dataset containing images and biomass measurements for model development, while predictions are evaluated on a separate hidden test set whose labels are not publicly available. 

The hidden test set is accessible only through the Kaggle evaluation platform. Participants submit predicted biomass values for the hidden images, and the platform returns only the resulting performance score, without revealing the ground-truth labels. Consequently, the hidden test set acts as an independent benchmark for assessing model generalization and reduces the risk of overfitting the evaluation procedure to the publicly available data. Final rankings are based on a weighted coefficient of determination ($R^2_w$) computed on this hidden data.

The author participated in this competition approximately seven months prior to writing this dissertation. During the competition, a notable discrepancy was observed between locally estimated performance and hidden-test performance. The public dataset consists exclusively of images collected in 2015, whereas the hidden test set contains images acquired during different sampling periods. Consequently, the hidden evaluation introduces temporal distribution shifts that are difficult to reproduce using standard validation procedures. The exact composition of the hidden test set is not publicly available.

These observations motivated a more systematic investigation of validation protocols and their ability to estimate real-world generalization performance. The resulting study led to a conference paper, currently under review, and to the present dissertation, which extends that work through additional analyses and discussion.

To achieve this objective, the following specific goals were established:

\begin{itemize}
    \item analysis of the characteristics of the CSIRO Image2Biomass dataset;
    \item implementation of a biomass estimation pipeline based on visual representations extracted with DINOv3 and regression models;
    \item evaluation of several validation strategies, including Random Cross-Validation, Date Cross-Validation, and Date-State Cross-Validation;
    \item comparison of the performance obtained in local validation with the performance obtained on the hidden test set of the Kaggle competition;
    \item quantification of the difference between locally estimated performance and actual performance on unseen data;
    \item investigation of the effects of temporal shifts on the generalization ability of the models.
\end{itemize}

\section{Research Methodology and Contributions}

For conducting the study, the CSIRO Image2Biomass dataset~\cite{liao2025estimating} was used, which contains RGB images of pasture plots together with measurements of several biomass components.

All experiments were carried out using the same preprocessing method, the same feature extractor, and the same predictive models, with only the data splitting strategy being modified. In this way, the observed differences can be directly attributed to the validation protocol and not to other experimental factors.

Visual features were extracted using the DINOv3 ViT-L/16 model in frozen mode~\cite{simeoni2025dinov3}. DINOv3 builds upon the self-supervised DINO family of vision transformers introduced by Caron et al.~\cite{caron2021dino}. Biomass prediction was then performed using two regression models: a Ridge model and a Multi-Layer Perceptron (MLP) neural network.

Performance evaluation was conducted using the weighted coefficient of determination $R^2_w$, the official metric of the Kaggle competition, which aggregates the $R^2$ values of the five biomass targets according to predefined weights. In addition, RMSE and MAE metrics were computed. For analyzing the reliability of the estimates, the optimism gap was defined and calculated as the difference between the performance estimated through local validation and the performance observed on the hidden test set.

The implementation was carried out in Python using several established libraries from the machine learning ecosystem: PyTorch for developing and training neural models, Scikit-Learn for regression and validation strategies, NumPy and Pandas for data processing, and Matplotlib and Seaborn for visualization. The hidden test set was evaluated through the Kaggle platform.

The main contributions of this work are:

\begin{itemize}
    \item conducting an experimental comparison between several validation protocols applied to a biomass estimation problem;
    \item quantifying the degree of optimism of different validation strategies by referencing the performance obtained on an independent and hidden test set;
    \item analyzing the impact of temporal and spatial shifts on model performance;
    \item developing an experimental pipeline for the systematic evaluation of validation strategies and the generalizability of machine learning models.
\end{itemize}

The implementation process made use of modern AI-assisted development tools for programming support, code refactoring, debugging assistance, and language refinement. However, the experimental design, validation methodology, model selection, result interpretation, and conclusions were determined and critically evaluated by the author.

\section{Relation to the Conference Paper}

This dissertation is based on a paper submitted to the \textit{International Conference on System Theory, Control and Computing (ICSTCC 2026)}, currently under review. The conference paper presents the core experimental comparison of validation protocols using the CSIRO Image2Biomass benchmark. The present dissertation extends that work by providing a more detailed theoretical background, a broader literature review and an expanded discussion of limitations and future directions. All experimental results presented here are consistent with those reported in the conference submission, with the dissertation offering a more comprehensive treatment of the methodology and findings.

\section{Thesis Structure}

The thesis is organized into five chapters.

The first chapter presents the research context, the motivation for choosing the topic, its relevance, the pursued objectives, the methodology used, and the main contributions of the work.

The second chapter provides a review of the specialized literature regarding biomass estimation, the use of machine learning and deep learning techniques in agriculture, as well as the validation methods used for evaluating model performance.

The third chapter describes the proposed experimental solution: the CSIRO Image2Biomass dataset, the preprocessing and feature-extraction pipeline, the Ridge and MLP regression models, the validation protocols, and the evaluation metrics.

The fourth chapter presents the experimental results, including local-validation scores, hidden-test scores, optimism gaps, per-target analyses, and the Leave-One-Period-Out temporal analysis.

The final chapter synthesizes the conclusions, highlights the contributions, discusses limitations, and proposes future research directions.

% \section{Context și motivație}

% Estimarea biomasei reprezintă o componentă importantă a agriculturii de precizie, fiind utilizată pentru monitorizarea stării culturilor și a pășunilor, evaluarea productivității și optimizarea procesului de gestionare a resurselor agricole. Metodele tradiționale de măsurare a biomasei se bazează în principal pe recoltarea și cântărirea materialului vegetal, proceduri care sunt costisitoare, consumatoare de timp și dificil de aplicat la scară largă.

% Progresele recente din domeniul inteligenței artificiale, al viziunii artificiale și al învățării automate au deschis posibilitatea estimării biomasei direct din imagini. Dezvoltarea modelelor bazate pe rețele neuronale profunde și utilizarea reprezentărilor vizuale pre-antrenate au permis obținerea unor rezultate promițătoare în numeroase aplicații agricole. Dintre diferitele surse de date disponibile, imaginile RGB sunt deosebit de atractive datorită costului redus, accesibilității și ușurinței de achiziție în condiții reale de exploatare.

% Cu toate acestea, performanțele ridicate raportate în literatura de specialitate sunt evaluate, în majoritatea cazurilor, folosind proceduri standard de validare, precum validarea încrucișată aleatorie. În prezența dependențelor temporale și spațiale specifice datelor agricole, aceste proceduri pot conduce la estimări excesiv de optimiste ale performanței reale. Astfel, un model care pare performant în timpul validării poate înregistra rezultate semnificativ mai slabe atunci când este aplicat pe date colectate în alte perioade sau în alte locații.

% Din această perspectivă, problema nu mai este doar construirea unui model capabil să estimeze biomasa, ci și evaluarea corectă a capacității acestuia de generalizare. În lipsa unor protocoale de validare adecvate, performanțele raportate pot oferi o imagine eronată asupra comportamentului modelului în condiții reale de utilizare.

% \section{Relevanța temei}

% Creșterea interesului pentru utilizarea inteligenței artificiale în agricultură a determinat apariția unui număr tot mai mare de modele dedicate estimării biomasei. Cu toate acestea, majoritatea cercetărilor se concentrează asupra îmbunătățirii performanței predictive, acordând o atenție redusă modului în care aceasta este evaluată.

% În practică, seturile de date agricole sunt caracterizate prin variații temporale și spațiale importante. Condițiile meteorologice, caracteristicile solului, tipurile de vegetație și momentul colectării imaginilor pot introduce schimbări de distribuție între datele utilizate pentru antrenare și cele întâlnite ulterior în exploatare. În aceste situații, metodele de validare convenționale pot să nu ofere o estimare realistă a performanței viitoare.

% Prin urmare, studiul fiabilității protocoalelor de validare reprezintă o problemă relevantă atât din perspectivă academică, cât și practică. Rezultatele unei astfel de analize pot contribui la alegerea unor proceduri de evaluare mai robuste și la o interpretare mai corectă a performanțelor raportate de modelele de învățare automată utilizate în domeniul agricol.

% \section{Obiectivele lucrării}

% Obiectivul principal al acestei lucrări este analizarea impactului diferitelor protocoale de validare asupra estimării performanței modelelor de predicție a biomasei în condiții de schimbare temporală și spațială a distribuției datelor.

% Pentru atingerea acestui obiectiv au fost stabilite următoarele obiective specifice:

%     \begin{itemize}
%     \item analiza caracteristicilor setului de date CSIRO Image2Biomass.
%     \item implementarea unui pipeline de estimare a biomasei bazat pe reprezentări vizuale extrase cu DINOv3 și modele de regresie.
%     \item evaluarea mai multor strategii de validare, incluzând Random Cross-Validation, Date Cross-Validation și Date-State Cross-Validation.
%     \item compararea performanțelor obținute în validarea locală cu performanțele obținute pe setul de test ascuns al competiției Kaggle.
%     \item cuantificarea diferenței dintre performanța estimată local și performanța reală pe date nevăzute.
%     \item investigarea efectelor schimbărilor temporale asupra capacității de generalizare a modelelor.
%     \end{itemize}

% \section{Metodologia cercetării}

% Pentru realizarea studiului a fost utilizat setul de date CSIRO Image2Biomass\cite{liao2025estimating}, care conține imagini RGB ale unor parcele de pășune împreună cu măsurători ale mai multor componente de biomasă.

% Toate experimentele au fost realizate utilizând aceeași metodă de preprocesare, același extractor de caracteristici și aceleași modele predictive, fiind modificată exclusiv strategia de împărțire a datelor. În acest mod, diferențele observate pot fi atribuite direct protocolului de validare și nu altor factori experimentali.

% Caracteristicile vizuale au fost extrase cu ajutorul modelului DINOv3 ViT-L/16 utilizat în regim frozen, iar predicția biomasei a fost realizată cu ajutorul a două modele de regresie: un model Ridge și o rețea neuronală de tip Multi-Layer Perceptron (MLP).

% Evaluarea performanței a fost realizată folosind indicatorul ($R^2_w$), metrică utilizată și în cadrul competiției Kaggle, precum și metricile RMSE și MAE. În plus, pentru analiza fiabilității estimărilor a fost definit și calculat optimism gap-ul, reprezentând diferența dintre performanța estimată prin validare și performanța observată pe setul de test ascuns.

% \section{Tehnologii și instrumente utilizate}

% Implementarea experimentelor a fost realizată în limbajul Python, utilizând mai multe biblioteci consacrate din ecosistemul machine learning:

%     \begin{itemize}
%     \item PyTorch pentru dezvoltarea și antrenarea modelelor neuronale.
%     \item Scikit-Learn pentru implementarea modelelor de regresie și a strategiilor de validare.
%     \item NumPy și Pandas pentru procesarea datelor.
%     \item Matplotlib și Seaborn pentru analiza și vizualizarea rezultatelor.
%     \item Kaggle pentru evaluarea performanței modelelor pe setul de test ascuns.
%     \end{itemize}
    
% De asemenea, pentru extragerea reprezentărilor vizuale a fost utilizat modelul DINOv3 ViT-L/16 preantrenat.

% \section{Contribuțiile lucrării}

% Principalele contribuții ale acestei lucrări sunt:

%     \begin{itemize}
%     \item realizarea unei comparații experimentale între mai multe protocoale de validare aplicate unei probleme de estimare a biomasei.
%     \item cuantificarea gradului de optimism al diferitelor strategii de validare prin raportare la performanța obținută pe un set de test independent și ascuns.
%     \item analiza impactului schimbărilor temporale și spațiale asupra performanței modelelor.
%     \item dezvoltarea unui pipeline experimental pentru evaluarea sistematică a strategiilor de validare și a capacității de generalizare a modelelor de învățare automată.
%     \end{itemize}

% \section{Structura lucrării}

% Lucrarea este organizată în patru capitole.

% Primul capitol prezintă contextul cercetării, motivația alegerii temei, relevanța acesteia, obiectivele urmărite, metodologia utilizată și contribuțiile principale ale lucrării.

% Capitolul al doilea realizează o analiză a literaturii de specialitate privind estimarea biomasei, utilizarea tehnicilor de învățare automată și învățare profundă în agricultură, precum și metodele de validare utilizate pentru evaluarea performanței modelelor.

% Capitolul al treilea descrie setul de date utilizat, arhitectura pipeline-ului implementat, protocoalele de validare analizate și rezultatele experimentale obținute. Sunt prezentate atât performanțele locale, cât și comparația cu performanțele raportate pe setul de test ascuns, împreună cu o analiză detaliată a erorilor și a efectelor schimbărilor temporale.

% Ultimul capitol sintetizează concluziile cercetării, evidențiază principalele contribuții ale lucrării, discută limitările existente și propune direcții pentru cercetări viitoare.